\begingroup
\renewcommand{\arraystretch}{1.45}
\fancyhead[L]{\small THEORY OF THE NEW TRANSCENDENTAL FUNCTIONS. III.}

\Needspace{27\baselineskip}
\textit{For division into five parts}\quad
$\left\{\begin{array}{cc}a&b\\A&B\end{array}\right\}$
\[
\begin{gathered}
\Bigl(\frac{a-A}{5A-a}\Bigr)^4=\frac{AA-BB}{aa-bb}\cdot\frac{BB}{bb}\\[8pt]
\Bigl(\frac{B-b}{5B-b}\Bigr)^4=\frac{AA-BB}{aa-bb}\cdot\frac{AA}{aa}\\[8pt]
\frac{a-A}{B-b}=\frac{5B-b}{5A-a}=\mathord{{}^{4}\!\surd}{\frac{aB}{bA}}
\end{gathered}
\]

\Needspace{14\baselineskip}
\[
\begin{gathered}
(a-A)(5A-a)^5=256Aabb(aa-bb)\\[8pt]
(B-b)(5B-b)^5=256Baab(aa-bb)
\end{gathered}
\]

\Needspace{14\baselineskip}
\[
\begin{gathered}
(a-A)^5(5A-a)=256AaBB(AA-BB)\\[8pt]
(B-b)^5(5B-b)=256AAbB(AA-BB)
\end{gathered}
\]

\Needspace{24\baselineskip}
\[
\begin{gathered}
a-A=\frac{4x[xx]^2[x^5][x^{20}]}{[x][x^4]}
 \\[8pt]\quad{}=2^{\frac43}\frac{a^{\frac16}A^{\frac16}B^{\frac5{12}}(AA-BB)^{\frac5{24}}}
 {b^{\frac1{12}}(aa-bb)^{\frac1{24}}}\\[1ex]
5A-a=4\frac{[x][x^4][x^{10}]^2}{[x^5][x^{20}]}
 \\[8pt]\quad{}=2^{\frac43}\frac{a^{\frac16}A^{\frac16}b^{\frac5{12}}(aa-bb)^{\frac5{24}}}
 {B^{\frac1{12}}(AA-BB)^{\frac1{24}}}
\end{gathered}
\]

\SourceLeaf{466}


\Needspace{27\baselineskip}
The following theorems belong to the theory of division into five parts. We denote
$1+xy\,.\,1+x^3y\,.\,1+x^5y\mathinner{.\,.\,.}1+\frac xy\,.\,1+\frac{x^3}y\,.\,1+\frac{x^5}y\mathinner{.\,.\,.}$
by $(x,y)$; then
\[
\begin{gathered}
[69]\quad{}
(x,\alpha y).(x,\tfrac ya)
\\[8pt]\quad{}=P\{1+xx(yy+\tfrac1{yy})+x^8(y^4+\tfrac1{y^4})+\mathinner{.\,.}\}\\[8pt]
\quad+Q\{y+\tfrac1y+x^4(y^3+\tfrac1{y^3})+x^{12}(y^5+\tfrac1{y^5})+\mathinner{.\,.}\}
\end{gathered}
\]
where $P$, $Q$ are independent of $y$.

\Needspace{13\baselineskip}
Thus
\[
\begin{gathered}(x,\alpha i)(x,\tfrac ai)=(xx,\alpha\alpha)
 \\[8pt]\quad{}=P(1-2xx+2x^8-\mathinner{.\,.})=\frac{[xx]^2}{[x^4]}P\end{gathered}
\]
\Needspace{9\baselineskip}
or
\[
P=(xx,\alpha\alpha)\frac{[x^4]}{[xx]^2}
\]

\Needspace{20\baselineskip}
Further, for $y=-\alpha x$
\[
\begin{gathered}
(x,-\alpha\alpha x).(x,-x)=0\\[8pt]
=P(1+\tfrac1{aa}+\alpha\alpha x^4+\tfrac{x^4}{a^4}+\mathinner{.\,.})
 \\[8pt]\quad{}-Q(\alpha x+\tfrac1{ax}+\tfrac x{a^3}+\alpha^3x^7+\tfrac{x^7}{a^5}+\mathinner{.\,.\,.})
\end{gathered}
\]
\Needspace{16\baselineskip}
that is
\[
\begin{gathered}
P\{\alpha+\tfrac1a+x^4(\alpha^3+\tfrac1{a^3})+x^{12}(\alpha^5+\tfrac1{a^5})+\mathinner{.\,.}\}\\[8pt]
\quad=\frac Qx\{1+xx(\alpha\alpha+\tfrac1{aa})+x^8(\alpha^4+\tfrac1{a^4})+\mathinner{.\,.}\}
\end{gathered}
\]
\Needspace{10\baselineskip}
Now
\[
P=\frac1{[xx]^2}\{1+xx(\alpha\alpha+\tfrac1{aa})+\mathinner{.\,.}\}
\]
\Needspace{10\baselineskip}
Thus
\[
Q=\frac x{[xx]}\{(\alpha+\tfrac1a)+x^4(\alpha^3+\tfrac1{a^3})+\mathinner{.\,.\,.}\}
\]

\Needspace{30\baselineskip}
and our theorem
\[
\begin{gathered}
70.\quad{}
(x,\alpha y).(x,\tfrac ya)\\[8pt]
=\frac1{[xx]^2}\left\{\begin{gathered}
\phantom{{}+x}\{1+xx(\alpha\alpha+\tfrac1{aa})+x^8(\alpha^4+\tfrac1{a^4})+\mathinner{.\,.}\}\\[8pt]
\quad\times\{1+xx(yy+\tfrac1{yy})+x^8(y^4+\tfrac1{y^4})+\mathinner{.\,.}\}\\[1ex]
+x\{\alpha+\tfrac1a+x^4(\alpha^3+\tfrac1{a^3})+x^{12}(\alpha^5+\tfrac1{a^5})+\mathinner{.\,.}\}\\[8pt]
\quad\times\{y+\tfrac1y+x^4(y^3+\tfrac1{y^3})+x^{12}(y^5+\tfrac1{y^5})+\mathinner{.\,.}\}
\end{gathered}\right\}
\end{gathered}
\]


\SourceLeaf{467}

\Needspace{18\baselineskip}
or
\[
\begin{gathered}(x,\alpha y).(x,\tfrac ya)
 =\frac{[x^4]^2}{[xx]^2}\{(xx,\alpha\alpha).(xx,yy)
 \\[8pt]\quad{}+x\alpha y(xx,\alpha\alpha xx).(xx,xxyy)\}\end{gathered}
\]

\Needspace{19\baselineskip}
(One can also easily make the series by which $P$ is multiplied $=0$,
by $y=ix$)
\[
71.\quad{}
(x,y)+(x,\tfrac xy)\mathord{{}^{4}\!\surd}{\frac x{yy}}
 =(x^{\frac14},y^{\frac12})\frac{[x^{\frac12}]}{[xx]}
\]

\Needspace{15\baselineskip}
\[
\begin{gathered}(x,\alpha x)=\frac1a(x,\tfrac xa),
\\[8pt]\quad{}\qquad
(x,\alpha xx)=\frac1{ax}(x,\alpha)\end{gathered}
\]

\Needspace{18\baselineskip}
Theorem 70 can also be stated as follows
\[
\begin{gathered}72.\quad{}
(x,\alpha).(x,\beta)=\frac{[x^4]^2}{[xx]^2}
 \{(xx,\alpha\beta).(xx,\tfrac a\beta)
 \\[8pt]\quad{}+x\alpha(xx,\alpha\beta xx).(xx,\tfrac{axx}\beta)\}\end{gathered}
\]

\Needspace{18\baselineskip}
From this it follows that
\[
\begin{gathered}(x,\tfrac xa).(x,\tfrac x\beta)=\frac{[x^4]^2}{[xx]^2}
 \{(xx,\tfrac{xx}{a\beta}).(xx,\tfrac a\beta)
 \\[8pt]\quad{}+\alpha(xx,\alpha\beta).(xx,\tfrac{axx}\beta)\}\end{gathered}
\]

\Needspace{19\baselineskip}
and further from this
\[
\begin{gathered}73.\quad{}
(x,\alpha).(x,\beta)+(x,\tfrac xa).(x,\tfrac x\beta)\mathord{\surd}{\frac x{a\beta}}
 \\[8pt]\quad{}=\frac{[x]^2}{[xx]^2}.(x^{\frac12},\mathord{\surd}{\alpha\beta}).(x^{\frac12},\mathord{\surd}{\tfrac\beta a})\end{gathered}
\]
\Needspace{25\baselineskip}
Now
\[
\begin{gathered}
\frac{1+2x+2x^4+\mathinner{.\,.\,.}}{1+2x^5+2x^{20}+\mathinner{.\,.\,.}}
\\[8pt]\quad{}=\frac{(ix^{\frac52},ix^{\frac32}).(ix^{\frac52},-ix^{\frac12})}
 {(ix^{\frac52},-ix^{\frac32}).(ix^{\frac52},ix^{\frac12})}\\[1ex]
\\[8pt]\quad{}=\frac{(-x^5,xx)(-x^5,-x)+x(-x^5,-x^3)(-x^5,x^4)}
 {(-x^5,xx)(-x^5,-x)-x(-x^5,-x^3)(-x^5,x^4)}
\end{gathered}
\]

\Needspace{28\baselineskip}
From this the first theorem to be proved follows directly. Similarly,
\[
\begin{gathered}
\frac{1+2x+2x^4+\mathinner{.\,.\,.}}{1+2x^5+2x^{20}+\mathinner{.\,.\,.}}\\[8pt]
\quad=\dfrac{\begin{gathered}1-\varepsilon x\,.\,1-\varepsilon\varepsilon x\,.\,1-\varepsilon^3x\,.\,1-\varepsilon^4x\,.\,\\[6pt]1+\varepsilon xx\,.\,1+\varepsilon\varepsilon xx\,.\,1+\varepsilon^3xx\,.\,1+\varepsilon^4xx\mathinner{.\,.\,.}\end{gathered}}
 {\begin{gathered}1+\varepsilon x\,.\,1+\varepsilon\varepsilon x\,.\,1+\varepsilon^3x\,.\,1+\varepsilon^4x\,.\,\\[6pt]1-\varepsilon xx\,.\,1-\varepsilon\varepsilon xx\,.\,1-\varepsilon^3xx\,.\,1-\varepsilon^4xx\mathinner{.\,.\,.}\end{gathered}}\\[1ex]
\quad=\frac{1-\varepsilon^4\,.\,1-\varepsilon^3}{1+\varepsilon^4\,.\,1+\varepsilon}
 \cdot\frac{(ix^{\frac12},i\varepsilon x^{\frac12}).(ix^{\frac12},i\varepsilon\varepsilon x^{\frac12})}
 {(ix^{\frac12},-i\varepsilon x^{\frac12}).(ix^{\frac12},i\varepsilon\varepsilon x^{\frac12})}\\[1ex]
\quad=\frac{\varepsilon\varepsilon-\varepsilon^3\,.\,\varepsilon-\varepsilon^4}
 {\varepsilon\varepsilon+\varepsilon^3\,.\,\varepsilon+\varepsilon^4}
 \cdot\frac{(-x,-\varepsilon^3x).(-x,\varepsilon)-\varepsilon x(-x,\varepsilon^3xx)(-x,-\varepsilon x)}
 {(-x,-\varepsilon^3x).(-x,\varepsilon)+\varepsilon x(-x,\varepsilon^3xx)(-x,-\varepsilon x)}\\[1ex]
\quad=\frac{-\varepsilon+\varepsilon\varepsilon+\varepsilon^3-\varepsilon^4}
 {+\varepsilon+\varepsilon\varepsilon+\varepsilon^3+\varepsilon^4}
 \cdot\frac{(x,\varepsilon^3x).(x,-\varepsilon)+\varepsilon^3(x,-\varepsilon^3)(x,\varepsilon x)}
 {(x,\varepsilon^3x).(x,-\varepsilon)-\varepsilon^3(x,-\varepsilon^3)(x,\varepsilon x)}
\end{gathered}
\]
From this the second theorem to be proved follows directly.

\clearpage
\endgroup
